\subsection{محدودیت‌ها و تحلیل شکست}
تحلیل کارکرد الگوریتم \lr{AdaSVD} موارد زیر را روشن می‌سازد:
\begin{itemize}
\item \textbf{پایداری قابل توجه در مدل‌های کدگذار:} دقت بالا در \lr{RoBERTa} (\lr{87.84\%}) و \lr{DistilBERT} (\lr{84.86\%}) اثربخشی مستقیم جذب شیفت میانگین و برازش \lr{ALS} را در حفظ فضای بازنمایی این مدل‌ها نشان می‌دهد.
\item \textbf{عملکرد قابل قبول در GPT-2:} ثبت دقت \lr{77.06\%} و F1 معادل \lr{0.8051} در مدل خودهمبسته GPT-2 با اشغال تنها \lr{393.51 MB} حافظه VRAM نشان‌دهنده استحکام فشرده‌سازی در لایه‌های پیش‌بینی کلمه بعدی است.
\item \textbf{افت دقت در Qwen (تحلیل رتبه موثر):} در مدل مقیاس بزرگ‌تر Qwen2.5، دستیابی به دقت \lr{50.90\%} ناشی از برش خشن ابعاد در لایه‌های \lr{gate\_proj} و \lr{up\_proj} بدون تنظیم دقیق ترمیمی بوده است. همچنین در سناریوهایی که اندازه دسته کالیبراسیون کمتر از ابعاد پنهان شبکه باشد، ماتریس کوواریانس ورودی دچار تکینگی شده و نیازمند تقویت منظم‌سازی تیکونوف در معکوس‌سازی ماتریس‌ها است.
\end{itemize}
